\section{Architecture}
\label{sec:architecture}
% 4 questions exist that need to be answered:
%   1. How do robots usually split the system --- from frame capture to motor actuation?
%   2. What options were considered?
%   3. What was the final decision? Why did we opt for it?
%   4. What are the tradeoffs?

\subsection{Structure of an Autonomous Robot}
An autonomous robot has two main goals: complete its task without human
intervention and perform it efficiently. Efficiency as a metric depends
entirely on what we are designing for --- an autonomous robot that can navigate
a track on its own. For the case of Navilott, this means a robot that can
navigate a street-style course on its own --- following a lane, stopping at
stop signs and obeying traffic lights.

In robotics, this is often called "autonomous navigation", a part of
"autonomous control," where the robot controls itself to perform a task, such
as following a line or avoiding obstacles~\citep{douglas2020autonav1}. A robot
can have many different methods to measure and sense its environment ---
camera, sonar, ultrasonic, IR, etc. Subprocesses in the robots software can use
these measurements to form a map of the environment. The robot's estimation
system can use the map to calculate the robot's position and orientation in the
world. The robot's navigation system can then use these estimates to determine
the robot's decision about what to do next. The robots control system can then
actuate motors in response to these decisions. Actuation is performed in two
layers: the control loop which sends PWM signals, and the driver which switches
battery voltages based on the PWM signals that actuate the motors. Feedback
from encoders on the motor can then be routed back into estimation and
navigation systems to ensure that the decision was taken in the right place.

The robot can then be divided into 6 main areas:

\begin{itemize}
    \item \textbf{Capture}: the robot's sensors and how they capture the environment
    \item \textbf{Perception}: the robot's processing algorithms and how they interpret those measurements
    \item \textbf{Estimation}: the robot's estimation system and how it generalizes those measurements
    \item \textbf{Navigation}: the robot's navigation system and how it forms decisions about the environment based on those perceptions
    \item \textbf{Control}: the robot's control system and how it moves the robot around those decisions
    \item \textbf{Actuation}: the robot's actuation system and how it actually moves the robot
\end{itemize}

\subsection{Single-Board Design}
% Think about: the criteria you judged by (control rate, $100 BOM, 12 V,
%   local processing, dev time); MCU-only vs Pi+MCU vs Pi-only; the PRD
%   recommends a microcontroller, so Pi-only departs from its guidance.
% Evidence: vision_stack/docs/architecture.md ("Decision: Pi Zero 2 W only")
%           MCU option and packet format: docs/archD1/
Our first concern was how the start of the stack would interact with the end of
it --- meaning how we would connect the robot's perception system to the motor
actuation. To achieve this, three options were considered:

\begin{itemize}
    \item \textbf{MCU-only:} An STM32 or Arduino runs the perception pipeline and
          drives the motors directly.
    \item \textbf{Pi-only:} The Pi runs the perception pipeline and drives the
          motors directly via its onboard GPIO pins.
    \item \textbf{Pi + MCU:} The Pi processes perception and sends a packet over
          UART to an STM32 or Arduino, which executes a fast, real-time PID loop
          and generates the PWM signals.
\end{itemize}

The Raspberry Pi was a necessary addition because of our robot's camera-centric
design. Intense computer vision needs a lot of RAM to run efficiently, ISP and
CSI-2 interface to capture video quickly and bypass CPU overhead, and OpenCV
library to perform the intense image processing we would require. For instance,
to process a frame that is $480\times270$ in color pixels would require 389 KB
of RAM --- more than what is found in most microcontrollers.

Integrating an MCU into our design would require a new protocol to transfer the
data between the Pi and the MCU, synchronization of clocks between the two
devices (the Pi captures and stamps frames based on its clock while encoders
from the motors would work on the MCU's clock), and managing two separate
codebases. These different challenges, while making the architecture more true
to a real-world robot, would also make it more difficult to maintain.

As such, a single-board design was chosen --- which sacrifices speed of a
dedicated motor control loop for the simplicity and testability of the Pi, as
well as accepting the overhead from the Linux scheduler. However, these
tradeoffs were considered when developing the stack and are scattered through
Sections~\ref{sec:software_architecture},~\ref{sec:limitations},
and~\ref{sec:results}.

\subsection{Additional Hardware}
% Think about: where the camera choice lives (here or Capture?), including
%   the IMX219 -> IMX290 swap; which parts were Ana's or Ignacio's calls.
% Evidence: docs/architecture.md (Parts; Pin and bus map); appendices/hardware.tex
Mapping the baseline architecture for the design introduced several new avenues
to explore for integrating hardware components. For a robotic system, a camera
alone would not be enough to function as a robot. How does the robot stay
powered, how does it move, what else would it use to sense its environment, how
can we monitor its internal state, and how does the control system communicate
with the user? Taking in these questions, we decided to add the following
hardware components:

\begin{table}[H]
    \centering
    \caption{Hardware added around the Pi.}
    \label{tab:hardware}
    \small
    \begin{tabularx}{\textwidth}{@{}XXXX@{}}
        \toprule
        \textbf{Part}                       & \textbf{Role}                                                 & \textbf{Reason}                                                                  & \textbf{Area}                    \\
        \midrule
        \textbf{Innomaker's Cam-MIPI290RAW} & Main vision sensor for the robot.                             & Allows the robot to visually see its environment.                                & Capture                          \\
        \textbf{ADS1115}                    & Measures the battery voltage; Pi has no analog input.         & Allows for our team to report the battery status to the user.                    & Capture (Battery)                \\
        \textbf{MPU6050}                    & Measures the robot's angular rate and acceleration.           & Allows us to measure the robots yaw and integrate it into our estimation system. & Capture $\rightarrow$ Estimation \\
        \textbf{TB6612FNG}                  & The robot's motor driver; Pi GPIO can't supply motor current. & Allows for precise control of the robot's motors.                                & Actuation                        \\
        \textbf{2 N20 Motors}               & The robot's actuation system.                                 & Allows for the robot to move.                                                    & Actuation                        \\
        \textbf{Encoders}                   & Measures the robot's wheel speeds.                            & Allows us to estimate the robot's actual speed.                                  & Estimation                       \\
        \textbf{TM1637}                     & The robot's run-time interface.                               & Allows us to display information about the robot's state to the user.            & Interface                        \\
        \textbf{Start Button}               & The robot's start button.                                     & Allows us to physically start the robot.                                         & Interface                        \\
        \bottomrule
    \end{tabularx}
\end{table}

Originally, we were using a Raspberry Pi Cam V2 module as our camera, which
featured an IMX219 sensor. However, its field of view is cropped by the mode
selection of the Pi. This added a heavy blindspot of 22 cm to the robot's
vision and made it incapable of viewing both the lane lines and the traffic
light in one frame while at a stop line. Details can be found in
Section~\ref{sec:capture}. In response, we opted for another type of camera:
the Innomaker's Cam-MIPI290RAW, a wide-angle camera that is compatible with the
Raspberry Pi's CSI-2 interface. This allowed us to achieve a fuller field of
view and reduced our blindspot down to 10 cm, but at the cost of the lens
producing a distortion effect which would require software to correct
(Section~\ref{sec:capture}).

\subsection{Software Language}
% --- Software Language ---
\TODO{What programming language was used for the software? Why was this language chosen over others? When it came to
    integration, what were the tradeoffs?}
% Think about: Python as glue over OpenCV's C++; the bindings it gave you
%   (OpenCV, GStreamer/libcamera, pigpio); the GIL and interpreter cost,
%   and whether that shows up in P1.

\subsection{Tooling}
% --- Tooling ---
\TODO{What tools were used to develop the software? What was gained and lost by using these tools? How did the tools affect
    the architecture?}
% Think about: Make (setup.mk once per Pi, session.mk every boot,
%   vision_stack/Makefile as the command list), pytest, git/PRs, the
%   linkers and debug views, AI-assisted sessions.
% Architecture shaped by tooling: one place for stage order; replayable
%   frame sources; estimation imports no hardware; diagnostics in a
%   separate process.
% Evidence: docs/guides/make.md; git log for when each .mk was added

\subsection{Software Architecture}\label{sec:software_architecture}
% --- Software Architecture ---
\TODO{After deciding on the main architecture and hardware peripherals to use, how was the software architecture designed?
    What were some of the key decisions made and why? How did the software evolve over time and what changes, if any, were made
    to the architecture?}
% Think about: phases and data flow; threads and the 50 ms frame budget;
%   two or three turning points, each with its reason.
% Evidence: docs/archD1/ and archive/run_pipeline.py (the earlier design)
% Figures available: four_phases.png, pi_threading.png
